\documentclass[11pt,a4paper]{article}
\usepackage[margin=25mm]{geometry}
\usepackage{amsmath,amssymb,amsthm,hyperref}
\newtheorem{proposition}{Proposition}
\title{A locally toroidal polytope with a nonorientable vertex figure}
\author{Human Compression project}\date{}
\begin{document}\maketitle
Fernandes, Leemans and Ivi\'c Weiss~\cite{source}, pp.~18--19, mention
the polytope $\{\{4,4\}_{(4,0)},\{4,6\}_3\}$ and state that its
group has order $768$ and its vertex figure has nonorientable genus four.
They explicitly leave further study of this example aside. We verify
these assertions separately from their tables of orientably regular and
chiral examples. Here the subscript in $\{4,6\}_3$ is the Petrie length.

\begin{proposition}
Let $U$ be the group generated by involutions $a,b,c,d$ with
\[
 (ab)^4=(bc)^4=(cd)^6=(ac)^2=(ad)^2=(bd)^2=1,
 \qquad (abcb)^4=(bcd)^3=1.
\]
Then $U$, with this ordered generating set, is a string C-group of order
$768$. Its facets are $\{4,4\}_{(4,0)}$, and its vertex figures are
the nonorientable maps $\{4,6\}_3$ of genus four. It is universal for
these two prescribed local presentations.
\end{proposition}
\begin{proof}
The program \path{verify.g} forms the displayed finitely presented group
itself. GAP's exact finite group algorithms give $|U|=768$ and a bijective
homomorphism to a permutation group of degree $16$. For clarity, the
images of the four distinguished generators are
\begin{align*}
a={}&(1\,2)(3\,6)(4\,7)(5\,8)(9\,11)(10\,12)(13\,14)(15\,16),\\
b={}&(1\,3)(2\,6)(4\,5)(7\,8)(9\,13)(10\,14)(11\,15)(12\,16),\\
c={}&(1\,4)(2\,7)(3\,9)(5\,10)(6\,11)(8\,12)(13\,14)(15\,16),\\
d={}&(1\,5)(2\,8)(3\,4)(6\,7)(9\,12)(10\,11)(13\,16)(14\,15).
\end{align*}
This permutation image alone would give only a lower bound on the
universal group. The calculation of the finitely presented group supplies
the upper bound and checks bijectivity. For every pair of subsets
$I,J\subseteq\{a,b,c,d\}$, the program checks
\[
 \langle I\rangle\cap\langle J\rangle=\langle I\cap J\rangle.
\]
These are all $256$ pairs, including the empty and full subsets. The
adjacent product orders are exactly $4,4,6$. The standard coset
construction for string C-groups therefore gives a regular abstract
polytope with this group and marking.

The facet and vertex stabilizers $\langle a,b,c\rangle$ and
$\langle b,c,d\rangle$ have orders $128$ and $48$. Independently forming
their three-generator local presentations gives these same orders; the
maps sending local generators to the displayed distinguished generators
are checked bijective. Thus there are no additional local relations.
The first presentation is the square torus $\{4,4\}_{(4,0)}$: in the
affine square reflection group, $abcb$ is a unit translation, and its
fourth power and its conjugates generate the period lattice $4\mathbb Z^2$.
The quotient has order $8\cdot16=128$. The second is exactly the
Petrie quotient $\{4,6\}_3$; the product $bcd$ has order three.

In the vertex figure the numbers of vertices, edges and faces are
respectively
\[
 [\langle b,c,d\rangle:\langle c,d\rangle]=4,\qquad
 [\langle b,c,d\rangle:\langle b,d\rangle]=12,\qquad
 [\langle b,c,d\rangle:\langle b,c\rangle]=6.
\]
The Euler characteristic is therefore $4-12+6=-2$. The relation
$(bcd)^3=1$ is a closed walk of odd length nine in its flag graph.
Orientability would assign opposite orientations to adjacent flag
triangles and hence make this graph bipartite. The odd closed walk rules
this out. The map is a connected closed nonorientable surface, so its
genus is $2-(-2)=4$.

Finally, any regular polytope with these local presentations satisfies
the displayed defining relations. Its distinguished generators therefore
receive a surjective homomorphism from $U$. Since $U$ itself has the
required intersection property and local groups, it is the universal
object for these local presentations.
\end{proof}

The group order and genus are assertions of the cited source, not new
claims of discovery. The new work is their explicit verification in the
present proof system. The finite presentation, all subgroup intersections
and the local isomorphisms are required exact computations in GAP; this
note is not a complete Lean proof. It settles this isolated source
assertion and does not answer any of that paper's open questions.

\begin{thebibliography}{9}
\bibitem{source} M.~E.~Fernandes, D.~Leemans and A.~Ivi\'c Weiss,
\emph{Hexagonal extensions of toroidal maps and hypermaps},
the retained 21-page 2018 edition, pp.~18--19.
\end{thebibliography}
\end{document}
